\documentclass[12pt, a4paper]{article}

% --- Encodage et langues ---
\usepackage[utf8]{inputenc}
\usepackage[T1]{fontenc}
\usepackage[french]{babel}

% --- Mise en page et typographie ---
\usepackage{geometry}
\geometry{top=2.5cm, bottom=2.5cm, left=2.5cm, right=2.5cm}
\usepackage{microtype}

% --- Packages fonctionnels ---
\usepackage{enumitem}
\usepackage{titlesec}
\usepackage{xcolor}
\usepackage{tcolorbox}
\usepackage{amssymb}
\usepackage{marvosym}
\usepackage{hyperref}

% --- Configuration des hyperliens ---
\hypersetup{
    colorlinks=true,
    linkcolor=darkgray,
    urlcolor=blue,
    pdftitle={Guide du don à une association en Afrique},
    pdfauthor={Abd El Hakim ZOUAÏ (Cutopia)},
    pdfsubject={Actions nécessaires pour un don sécurisé et optimisé fiscalement}
}

% --- Couleurs personnalisées ---
\definecolor{lightgray}{gray}{0.95}
\definecolor{darkgray}{gray}{0.4}
\definecolor{darkgreen}{rgb}{0.0, 0.5, 0.0}
\definecolor{boxblue}{rgb}{0.85, 0.92, 0.95}
\definecolor{cutopia}{rgb}{0.85, 0.35, 0.55}
\definecolor{orangebox}{rgb}{0.98, 0.75, 0.5}
\definecolor{redalert}{rgb}{0.8, 0.2, 0.2}

% --- Style des sections ---
\titleformat{\section}
{\normalfont\Large\bfseries\color{darkgray}}
{\thesection}{1em}{}

\titleformat{\subsection}
{\normalfont\large\bfseries\color{darkgray}}
{\thesubsection}{1em}{}

% --- Environnements ---
\newtcolorbox{resumebox}{
    colback=lightgray,
    colframe=darkgray,
    arc=0mm,
    boxrule=0.5pt,
    boxsep=5pt,
    left=10pt, right=10pt, top=10pt, bottom=10pt
}

\newtcolorbox{helpbox}{
    colback=boxblue,
    colframe=blue!50!black,
    arc=2mm,
    boxrule=0.8pt,
    boxsep=5pt,
    left=10pt, right=10pt, top=8pt, bottom=8pt
}

\newtcolorbox{alertbox}{
    colback=redalert!10!white,
    colframe=redalert,
    arc=2mm,
    boxrule=1pt,
    boxsep=5pt,
    left=10pt, right=10pt, top=8pt, bottom=8pt,
    fonttitle=\bfseries,
    title={[!] Alerte}
}

\newtcolorbox{conseilbox}{
    colback=cutopia!8!white,
    colframe=cutopia,
    arc=3mm,
    boxrule=1.2pt,
    boxsep=6pt,
    left=12pt, right=12pt, top=10pt, bottom=10pt,
    fonttitle=\bfseries\large,
    title={[*] Conseil gratuit Cutopia [*]}
}

% --- Symboles simples sans fontawesome ---
\newcommand{\cmark}{\textcolor{darkgreen}{\textbf{\checkmark}}}
\newcommand{\arrow}{\textcolor{darkgreen}{\textbf{$\rightarrow$}}}
\newcommand{\heart}{\textcolor{cutopia}{\textbf{$\heartsuit$}}}
\newcommand{\starred}{\textcolor{darkgray}{\textbf{$\star$}}}
\newcommand{\euro}{\textcolor{darkgreen}{\textbf{\EUR}}}
\newcommand{\info}{\textcolor{blue}{\textbf{(i)}}}

% --- Listes personnalisées ---
\newlist{checklist}{itemize}{2}
\setlist[checklist]{label=\cmark, leftmargin=*}

\newlist{arrowlist}{itemize}{2}
\setlist[arrowlist]{label=\arrow, leftmargin=*}

\newlist{heartlist}{itemize}{2}
\setlist[heartlist]{label=\heart, leftmargin=*}

\newlist{starredlist}{itemize}{2}
\setlist[starredlist]{label=\starred, leftmargin=*}

% --- Début du document ---
\begin{document}

% --- En-tête ---
\begin{center}
    {\Huge \textbf{Guide complet du don sécurisé}} \\[0.3cm]
    {\Large \textbf{Vers une association en Afrique}} \\[0.3cm]
    {\large \textcolor{darkgray}{Pour les donneurs européens}} \\[0.5cm]
    \rule{0.7\textwidth}{0.4pt} \\[0.3cm]
    \textit{Un accompagnement étape par étape pour un don transparent, efficace et optimisé fiscalement}
\end{center}

\vspace{0.5cm}

% --- Table des matières ---
\tableofcontents
\newpage

% --- Introduction ---
\section*{Introduction}
\addcontentsline{toc}{section}{Introduction}
Ce guide a été conçu pour vous accompagner dans toutes les étapes de votre don à une association en Afrique. Que vous soyez un donateur occasionnel ou régulier, vous trouverez ici les informations nécessaires pour agir en toute sécurité et maximiser l'impact de votre générosité.

\vspace{0.3cm}
\begin{resumebox}
    \centering
    \textbf{\Large Ce que vous trouverez dans ce guide} \\[0.3cm]
    \begin{tabular}{ll}
        \cmark & Les vérifications préalables à tout don \\
        \euro & Les modes de transfert les plus adaptés \\
        \starred & La sécurisation de votre don \\
        \info & La défiscalisation (selon votre pays) \\
        \arrow & Les documents à exiger de l'association \\
        \heart & Une assistance gratuite via le réseau Cutopia
    \end{tabular}
\end{resumebox}
\vspace{0.3cm}

\newpage

% --- Section 1 ---
\section{Vérifications préalables obligatoires}

Avant tout transfert d'argent, vous devez impérativement vous assurer de la légitimité et de la fiabilité de l'association destinataire.

\subsection{Identifier clairement l'association}
\begin{checklist}
    \item Nom exact de l'association (tel qu'enregistré officiellement)
    \item Pays d'intervention (peut être différent du pays d'enregistrement)
    \item Objet social principal : santé, éducation, environnement, droits humains, etc.
    \item Année de création et ancienneté des dirigeants
\end{checklist}

\subsection{Vérifier la légalité dans son pays d'origine}
\begin{checklist}
    \item L'association est-elle enregistrée auprès des autorités compétentes ?
    \begin{itemize}
        \item En Afrique francophone : récépissé ministériel (Intérieur ou Administration territoriale)
        \item En Afrique anglophone : certificate of incorporation
        \item En Afrique lusophone : registro da associação
    \end{itemize}
    \item Possède-t-elle un numéro d'identification fiscale (NIF) ?
    \item Est-elle à jour dans ses déclarations annuelles ?
\end{checklist}

\subsection{Rechercher des avis et rapports de transparence}
\begin{checklist}
    \item Consulter le rapport annuel de l'association (au moins les deux derniers exercices)
    \item Vérifier sa présence sur des plateformes reconnues :
    \begin{itemize}
        \item GiveWell (recommandation basée sur l'efficacité démontrée)
        \item GuideStar (transparence financière)
        \item GlobalGiving (projets vérifiés)
    \end{itemize}
    \item Rechercher des articles de presse ou études de cas impliquant l'association
\end{checklist}

\subsection{Contacter directement l'association}
\begin{arrowlist}
    \item Envoyez un courriel ou appelez avant de donner
    \item Posez des questions précises
    \item Évaluez la qualité et la rapidité de leur réponse
\end{arrowlist}

\begin{helpbox}
    \centering
    \info\ \textbf{Bon à savoir} \\
    Une association fiable et bien structurée répondra dans un délai raisonnable (48 \`a 72 heures) et fournira sans difficulté les informations demandées.
\end{helpbox}

\newpage

% --- Section 2 ---
\section{Choix du mode de transfert}

\subsection{Comparatif des modes de transfert}

\begin{tcolorbox}[colback=lightgray, colframe=darkgray, title=Tableau comparatif]
    \small
    \renewcommand{\arraystretch}{1.5}
    \begin{tabular}{|p{3cm}|p{2cm}|p{2cm}|p{2.5cm}|p{2cm}|}
        \hline
        \textbf{Mode} & \textbf{Frais} & \textbf{Délai} & \textbf{Sécurité} & \textbf{Recommandation} \\
        \hline
        Virement SWIFT & 20-50 € + change & 2-5 jours & Élevée & Excellent \\
        \hline
        Wise (TransferWise) & 0.5-1\% & 1-2 jours & Élevée & Excellent \\
        \hline
        Revolut & 0-1\% & Instantané & Élevée & Très bon \\
        \hline
        Plateforme (GlobalGiving) & 5-15\% & Variable & Élevée & Bon \\
        \hline
        Mobile Money & 1-5\% + change & Instantané & Moyenne & Moyen \\
        \hline
        Western Union & 5-15\% & Instantané & Faible & À éviter \\
        \hline
        Espèces & 0\% & Instantané & Très faible & À éviter \\
        \hline
    \end{tabular}
\end{tcolorbox}

\subsection{Détails par mode de transfert}

\subsubsection{Virement bancaire international (SWIFT)}
\begin{arrowlist}
    \item \textbf{À demander à l'association :} IBAN, code BIC/SWIFT
    \item \textbf{Vérification :} Le nom du titulaire doit être celui de l'association
    \item \textbf{À savoir :} Les frais intermédiaires peuvent s'ajouter
\end{arrowlist}

\subsubsection{Virement SEPA}
\begin{arrowlist}
    \item \textbf{Condition :} Compte dans une banque de la zone SEPA
    \item \textbf{Avantage :} Frais généralement nuls
\end{arrowlist}

\subsubsection{Plateformes de don sécurisées}
\begin{arrowlist}
    \item GlobalGiving : projets vérifiés
    \item HelloAsso : pour les projets portés par des associations françaises
\end{arrowlist}

\begin{alertbox}
    \centering
    \textbf{Modes à éviter absolument} \\
    \vspace{0.2cm}
    Western Union, envoi d'espèces, compte bancaire personnel
\end{alertbox}

\newpage

% --- Section 3 ---
\section{Optimisation des frais et du change}

\subsection{Comparer les frais bancaires}
\begin{starredlist}
    \item Vérifiez la grille tarifaire de votre banque
    \item Attention aux frais cachés (commission fixe + frais de change)
\end{starredlist}

\subsection{Utiliser un service de transfert à bas coût}
\begin{starredlist}
    \item \textbf{Wise :} Recommandé, frais transparents (0.5-1\%)
    \item \textbf{Revolut :} Très compétitif
\end{starredlist}

\subsection{Choisir la devise d'envoi}
\begin{starredlist}
    \item Demandez à l'association quelle devise est la plus pratique
    \item Envoi en euros ou dollars : souvent avantageux
\end{starredlist}

\newpage

% --- Section 4 ---
\section{Sécurisation du don}

\begin{checklist}
    \item \textbf{Ne jamais donner sans reçu}
    \item \textbf{Exiger un document officiel}
    \item \textbf{Protection des données personnelles}
    \item \textbf{Signaler tout comportement suspect}
\end{checklist}

\subsection{Reçu fiscal et attestation}
\begin{arrowlist}
    \item L'association doit fournir une attestation avec :
    \begin{itemize}
        \item Son nom, adresse et numéro d'enregistrement
        \item Votre nom et adresse
        \item La date et le montant
    \end{itemize}
\end{arrowlist}

\newpage

% --- Section 5 ---
\section{Défiscalisation selon votre pays}

\subsection{France}
\begin{heartlist}
    \item Taux : 66\% du don (limite de 20\% du revenu)
    \item L'association doit avoir un relais fiscal en France
\end{heartlist}

\subsection{Belgique}
\begin{heartlist}
    \item Taux : 45\% du don (minimum 40 €)
    \item L'association doit être reconnue
\end{heartlist}

\subsection{Suisse}
\begin{heartlist}
    \item Taux variable selon les cantons (20-30\%)
\end{heartlist}

\newpage

% --- Section 6 ---
\section{Suivi après le don}

\subsection{Conservation des justificatifs}
\begin{checklist}
    \item Courriel de confirmation
    \item Relevé bancaire
    \item Attestation de don
\end{checklist}

\subsection{Demander un retour d'impact}
\begin{arrowlist}
    \item Envoyez une demande 3 à 6 mois après le don
    \item Demandez un rapport, des photos, des indicateurs
\end{arrowlist}

\newpage

% --- Section 7 ---
\section{Liste des documents à exiger}

\begin{tcolorbox}[colback=darkgreen!5!white, colframe=darkgreen!75!black, title=\textbf{Catégorie 1 : Preuve d'existence légale}]
\begin{starredlist}
    \item Statuts de l'association
    \item Récépissé de déclaration
    \item Numéro d'Identification Fiscale (NIF)
\end{starredlist}
\end{tcolorbox}

\begin{tcolorbox}[colback=blue!5!white, colframe=blue!75!black, title=\textbf{Catégorie 2 : Gouvernance et transparence}]
\begin{starredlist}
    \item Procès-verbal de la dernière AG
    \item Composition du Bureau
    \item Rapports d'activité
    \item États financiers annuels
\end{starredlist}
\end{tcolorbox}

\begin{tcolorbox}[colback=orange!5!white, colframe=orange!75!black, title=\textbf{Catégorie 3 : Documents du projet}]
\begin{starredlist}
    \item Dossier de présentation
    \item Budget prévisionnel
    \item Rapports d'exécution
\end{starredlist}
\end{tcolorbox}

\begin{tcolorbox}[colback=darkgray!5!white, colframe=darkgray!75!black, title=\textbf{Catégorie 4 : Documents bancaires}]
\begin{starredlist}
    \item RIB officiel au nom de l'association
    \item Lettre officielle de demande
    \item Convention de partenariat
\end{starredlist}
\end{tcolorbox}

\begin{alertbox}
    \centering
    \textbf{Signal d'alarme} \\
    Si l'association refuse de fournir ces documents, suspendez le transfert.
\end{alertbox}

\newpage

% --- Section 8 ---
\section*{Conseil gratuit Cutopia}
\addcontentsline{toc}{section}{Conseil gratuit Cutopia}

\begin{conseilbox}
    \centering
    {\Large \textbf{Abd El Hakim ZOUAÏ}} \\[0.2cm]
    \textit{Conseiller bénévole chez Cutopia} \\[0.3cm]
    
    Consultation gratuite pour :
    \begin{starredlist}
        \item Vérifier la légalité d'une association africaine
        \item Structurer un don transfrontalier sécurisé
        \item Obtenir un modèle de convention de partenariat
    \end{starredlist}
    
    \vspace{0.3cm}
    \rule{0.7\textwidth}{0.3pt} \\[0.3cm]
    \textbf{Contact :} \\
    \texttt{abdelhakimzouai@gmail.com} \\[0.2cm]
    \texttt{+213 672 41 01 92}
\end{conseilbox}

\newpage

% --- Checklist finale ---
\section*{Checklist finale}
\addcontentsline{toc}{section}{Checklist finale}

\begin{tcolorbox}[colback=lightgray, colframe=darkgray]
    \begin{checklist}
        \item[$\square$] J'ai vérifié l'existence légale de l'association
        \item[$\square$] J'ai analysé ses états financiers
        \item[$\square$] J'ai choisi un mode de transfert sécurisé
        \item[$\square$] J'ai obtenu un RIB au nom de l'association
        \item[$\square$] J'ai conservé tous les justificatifs
    \end{checklist}
\end{tcolorbox}

\vspace{1cm}
\begin{center}
    \heart\ \textcolor{cutopia}{\textbf{Votre don peut changer des vies s'il est bien fait.}} \ \heart
\end{center}

\vspace{0.5cm}
\begin{center}
    \rule{0.6\textwidth}{0.3pt} \\[0.3cm]
    {\small \textcolor{darkgray}{Guide r\'edig\'e avec le soutien d'\textbf{Abd El Hakim ZOUA\"I (Cutopia)} pour \textbf{Ethic Luxury}.}} \\
    \small{\textcolor{darkgray}{\texttt{abdelhakimzouai@gmail.com} \quad | \quad \texttt{+213 672 41 01 92}}}
\end{center}

\end{document}
